\documentclass[11pt,a4paper]{article}
\usepackage[margin=1.5cm]{geometry}
\usepackage{enumitem}
\usepackage[hidelinks]{hyperref}
\usepackage{titlesec}
\usepackage{xcolor}

\definecolor{accent}{HTML}{1F4FD8}

\titleformat{\section}{
	\large\bfseries\color{accent}
}{}{0em}{}[\titlerule]

\pagenumbering{gobble}

\begin{document}
	
	% ===== HEADER =====
	\begin{center}
		{\Huge \textbf{Ricardo Ramirez Caero}}\\
		\vspace{2mm}
		\href{mailto:ricardo@richixs.dev}{ricardo@richixs.dev} $\mid$
		\href{https://github.com/Richixs}{GitHub} $\mid$
		\href{https://linkedin.com/in/richixs}{LinkedIn} $\mid$
		Cochabamba, Bolivia
	\end{center}
	
	% ===== PROFILE =====
	\section*{Perfil Profesional}
	Desarrollador \textbf{Full Stack} y estudiante de Ingeniería Informática con experiencia como freelancer en el desarrollo de aplicaciones web, backend y mobile. He participado en proyectos reales para instituciones, ONGs y organizaciones, utilizando tecnologías modernas y buenas prácticas. Apasionado por el \textbf{open source}, Linux, la automatización y la participación en comunidades tecnológicas y hackatones.
	
	% ===== SKILLS =====
	\section*{Habilidades Técnicas}
	\begin{itemize}[leftmargin=*]
		\item \textbf{Frontend:} Angular, React, HTML, CSS, Tailwind
		\item \textbf{Backend:} Django, FastAPI, NestJS, Laravel, PHP, APIs REST, GraphQL
		\item \textbf{Mobile:} Android nativo (Java / Kotlin)
		\item \textbf{ERP:} Odoo (Python, PostgreSQL)
		\item \textbf{Bases de Datos:} PostgreSQL, MySQL, MongoDB
		\item \textbf{DevOps:} Docker, Kubernetes, ArgoCD, CI/CD
		\item \textbf{Otros:} Linux, Bash, Git, GitHub, GitLab
	\end{itemize}
	
	% ===== EXPERIENCE / PROJECTS =====
	\section*{Experiencia y Proyectos}
	\textbf{Desarrollador Freelancer} \hfill 2025 -- Presente
	\begin{itemize}[leftmargin=*]
		\item Desarrollo de una \textbf{aplicación en Python} para reconocimiento de \textbf{Equipo de Protección Personal (EPP)} mediante cámara, utilizando visión por computadora (\textbf{OpenCV}, \textbf{Python}).
		\item Desarrollo del sistema de gestión para el centro de diálisis \textbf{Hemosalud} utilizando \textbf{Angular} y \textbf{Django}.
		\item Desarrollo de \textbf{ChekiBot}, agente de IA para responder consultas electorales y combatir la desinformación para \textbf{Chequea Bolivia}, incluyendo diseño de APIs y despliegue backend (\textbf{React}, \textbf{FastAPI}).
		\item Desarrollo de \textbf{Gaceta}, plataforma web para la publicación de decretos y normativas gubernamentales usando \textbf{PHP} y \textbf{Laravel}.
		\item Desarrollo de un \textbf{sistema de cobros} para la Asociación de Agua \textbf{Willcataco} utilizando \textbf{Angular} y \textbf{FastAPI}.
	\end{itemize}
	
	% ===== EDUCATION =====
	\section*{Educación}
	\textbf{Ingeniería Informática} (En curso) \\
	Universidad Mayor de San Simón (UMSS) \hfill 2023 -- Actualidad
	
	% ===== COMMUNITY =====
	\section*{Comunidad y Open Source}
	\begin{itemize}[leftmargin=*]
		\item Miembro activo y organizador de \textbf{HackLab Brickheads} -- Cochabamba.
		\item Miembro activo de la \textbf{SCESI} (Sociedad Científica de Estudiantes de Sistemas e Informática) - UMSS.
				Expositor del taller \textit{``Linux a tu manera: introducción y personalización''}, enfocado en uso de Linux, personalización del entorno y herramientas básicas.
			
		\item Participación activa en eventos de la comunidad tecnológica y \textbf{hackatones}.
		\item Interés y contribución en proyectos \textbf{open source}.
	\end{itemize}
	
\end{document}
